% This is a composition with an existing incidence construction, not a new SAT order.
\section{A direct and explicit bridge to 3-SAT}
\label{sec:occurrence-bridge}

The occurrence construction below is already present in
\texttt{eth\_balanced\_packing/report.tex}, Section 8 (``Consequences
for SAT and limits of the improvement''), in the paragraph beginning
``For a concrete occurrence-based application.'' That source fixes
variables of degree greater than $\lceil(N+1)^{9/10}\rceil$.
We give the full composition for an arbitrary root-degree promise,
including a logarithmic threshold~\cite{balanced}. This avoids comparing a new
contractor only against the weaker all-right equality-tree encoding.

\begin{proposition}[Occurrence-preserving 3-SAT bridge]
\label{prop:occurrence-bridge}
Let $F$ be a 3-CNF, let $L$ be its number of literal occurrences after
removing tautologies and duplicate literals within clauses, and put
$S=L+1$. Suppose a uniform contraction algorithm on bipartite Boolean
networks with at most $S$ wires, root precision at most $S$, and
maximum degree at most $d(S)$ has time bound
\[
 (S+\mathcal L)^{O(1)}
 2^{(C+o(1))S/\sqrt{\log_2 S}},
\]
where $\mathcal L$ is the description length and the $o(1)$ is uniform
over admissible networks. Let an integer threshold $D(S)$ satisfy
\[
 \max\{3,D(S)\}\le d(S),\qquad
 D(S)/\sqrt{\log_2 S}\longrightarrow\infty.
\]
Then 3-SAT decision, and in fact exact counting, has time bound
\[
 \operatorname{poly}(n,m)\,
 2^{(C+o(1))L/\sqrt{\log_2 L}}.
\]
If the contractor uses polynomial space, so does this algorithm.
If its space coefficient is $C_S$ on the same normalized scale,
the composed space coefficient remains $C_S$.
\end{proposition}

\begin{proof}
An empty clause and a formula with no remaining clauses are handled
directly. Let $n_0\le L$ be the number of occurring variables. All other
declared variables contribute a known factor $2^{n-n_0}$ to the count
and do not affect the decision answer. For each occurring variable $v$,
let $d_v$ be its original occurrence degree, and set
\[
 H=\{v:d_v>D(S)\},\qquad h=|H|.
\]
Because the occurrences partition by variable,
$h(D(S)+1)\le L$. Enumerate all $2^h$ assignments to $H$.

For a branch, use the original variable--clause incidence graph after
removing just the ports belonging to variables in $H$. Each unfixed
variable carries the tensor that is one precisely when all its incident
bits agree. Each clause carries its satisfaction tensor after the fixed
values have been substituted. Retain every unfixed original incidence,
even when a fixed literal has already made a clause identically true;
the latter clause tensor is simply constant one on those ports. An
isolated clause is a scalar zero or one. Retaining these incidences
ensures that unfixed variables in already satisfied clauses are still
counted correctly.

The network is bipartite, has
\[
 M=L-\sum_{v\in H}d_v\le L
\]
wires, clause degree at most three, and variable degree at most $D(S)$.
Its equality and clause entries are computable in polynomial time.
Every assignment to the $n_0-h$ unfixed occurring variables determines
exactly one incidence assignment with nonzero equality tensors.
Hence the contraction is exactly the branch's number of satisfying
assignments.

Use the scalar ring $\mathbb Z/2^{P_0}\mathbb Z$ with
$P_0=n_0-h+1\le S$. The branch count lies in $[0,2^{n_0-h}]$,
so its least nonnegative residue determines it exactly. Sum these
counts in an $(n_0+1)$-bit accumulator, then multiply by the known
power of two for unused declared variables. The same procedure
decides satisfiability by testing whether the final count is positive.

The branching exponent obeys
\[
 h\le \frac{L}{D(S)+1}
 =o\!\left(\frac{S}{\sqrt{\log_2 S}}\right).
\]
All graphs and entry descriptions have size polynomial in the input.
Thus multiplying the contractor's time by $2^h$ leaves its leading
coefficient $C$ unchanged. Finally,
$S/\sqrt{\log_2 S}=(1+o(1))L/\sqrt{\log_2 L}$.
Branches are processed sequentially, so only polynomial additional
storage is needed. This proves both space assertions.
\end{proof}

Taking $D(S)=\lceil\log_2 S\rceil$ gives logarithmic root degree and
branching exponent $O(L/\log L)$, which is lower order than the
contraction exponent. In particular, the bridge meets every fixed
positive degree-power promise $d(S)=\lceil S^\alpha\rceil$ eventually.
Alternatively, any fixed $0<\delta<\alpha$ allows
$D(S)=\lfloor S^\delta\rfloor$. Any degree promise that eventually dominates
$\lceil\log_2 S\rceil$ can use the logarithmic-threshold version directly;
one must still verify the degrees of the fresh tensors it creates.

Because $L\le3m$, the occurrence bound implies
\[
 \operatorname{poly}(n,m)
 2^{(3C+o(1))m/\sqrt{\log_2 m}}.
\]
For fixed $D>0$ and $m\le Dn$, it also implies
\[
 \operatorname{poly}_D(n,m)
 2^{(3DC+o(1))n/\sqrt{\log_2 n}}.
\]
The latter statement follows from eventual monotonicity of
$x/\sqrt{\log_2 x}$ and $L\le3Dn$; unused declared variables create
no extra wire or precision term. These are explicit upper-bound
coefficients, not measured SAT speedups. At unrestricted density,
the occurrence route must be combined with the existing compressed
reduction; by itself it does not improve the worst-case $n$-only order.
